\begin{center}
	\rule{\textwidth}{1.5pt}
	\vspace{0.5pt}
	
	{\LARGE\bfseries Skyrmion Detection and Phase Classification}
	
	{\large P. Haas$^{[1]}$, I. Rasenberg$^{[2]}$}
	
	{\small$[1]$ Theoretical Physics, Utrecht University. $[2]$ Experimental Physics, Utrecht University.}
	
	{\small \today}
	
	\rule{\textwidth}{1.5pt}
\end{center}

	\begin{abstract}
		
		Skyrmions are topologically protected spin textures that emerge in magnetic systems with competing exchange, Dzyaloshinskii--Moriya (DM), and Zeeman interactions. In this work, a classical spin model on a two-dimensional lattice is investigated using Monte Carlo simulations to study the formation, detection, and thermal stability of skyrmions as a function of dimensionless DM interaction strength \(D/J\), applied magnetic field \(H/J\), and temperature.
		
		A discretized topological charge density is computed from spin configurations and used as the basis for a skyrmion detection algorithm. Local extrema in the smoothed charge density are identified and classified to distinguish isolated skyrmionic textures from spiral and ferromagnetic phases. This allows the construction of phase diagrams in the \((D/J, H/J)\) parameter space.
		
		The phase diagram reveals distinct regions corresponding to ferromagnetic order, spiral spin textures, and skyrmion formation. The skyrmionic region is further analysed as a function of temperature by studying both skyrmion formation upon heating and skyrmion stability upon thermal destruction. Cutoff temperatures for both processes are extracted, providing a quantitative measure of thermal stability across the phase diagram.
		
		\textbf{[conclusion here]}
		
	\end{abstract}
	
	\noindent\textbf{Keywords:}
	Skyrmions, Magnetic Phases, Topological Magnetism, Spin Textures, Phase Classification
	

